\paragraph{The regulator fixes the class of linear contours (algebraic at every \(N\), proved at \(N=2\)).}
Keep the setting and notation of the preceding paragraph. On the contour \(A_\mu=\sum_\nu\Lambda_{\mu\nu}X_\nu\) the regulator of Definition~(A) is \(\epsilon\sum_{\alpha,\beta}M_{\alpha\beta}\operatorname{Tr}X_\alpha X_\beta\) with \(M=\Lambda^{\mathrm T}\Lambda\). The matrix \((\operatorname{Tr}X_\alpha X_\beta)\) is a real Gram matrix, and every real positive semidefinite matrix of rank one occurs, so the regulator factor has modulus at most \(1\) on the whole contour exactly when \(\operatorname{Re}M\) is positive semidefinite.

(i) \emph{Inertia.} Let \(\ell=a+\mathrm ib\), with \(a,b\) real, be the row \(\mu=0\) of \(\Lambda\). Then \(h=M-2\ell\ell^{\mathrm T}\), so \(\operatorname{Re}h=\operatorname{Re}M-2aa^{\mathrm T}+2bb^{\mathrm T}\), and \(v=\Lambda^{-1}e_0\) satisfies \(hv=-Mv\), so \(\bar v^{\mathrm T}(\operatorname{Re}h)v=-\bar v^{\mathrm T}(\operatorname{Re}M)v\). If \(\operatorname{Re}M\) is positive semidefinite, \(\operatorname{Re}h\) therefore has at most one negative eigenvalue and is not positive definite. If \(\operatorname{Re}M\) is positive definite, \(\operatorname{Re}h\) has exactly one negative and nine positive eigenvalues, the inertia of \(\eta\). By (ii) of the preceding paragraph, \(\operatorname{Re}h\) has at least nine negative eigenvalues on \(-\mathfrak C_0\), and so at least nine non-positive ones on its closure. Hence \(h\notin-\mathfrak C_0\) in the first case, and \(h\) is not in the closure of \(-\mathfrak C_0\) in the second. In particular \(\operatorname{Re}h\) is never positive definite, whereas the Euclidean contours \(h=\mathrm e^{\mathrm i\pi/4}T^{\mathrm T}T\), \(T\) real, have \(\operatorname{Re}h\) positive definite. These statements are algebraic and hold for every \(N\). The same proof, in a real basis in which \(P\) is the identity and \(\eta\) is diagonal, gives them for the regulator \(\epsilon\sum_{\mu,\nu}P_{\mu\nu}\operatorname{Tr}A_\mu A_\nu\), with \(M=\Lambda^{\mathrm T}P\Lambda\), whenever \(P\) is positive definite and \(P^{-1}\eta\) has a nine-fold eigenvalue, for example \(P=c\,L^{\mathrm T}L\) with \(c>0\) and \(L\) a real Lorentz transformation. With the mass term the quadratic part of the exponent is \(-\sum_{\mu,\nu}(\epsilon\delta_{\mu\nu}+\mathrm i\gamma'\eta_{\mu\nu})\operatorname{Tr}A_\mu A_\nu\), where \(\gamma'=\gamma/2g_{\mathrm{YM}}^2\). Write \(M_\gamma=\epsilon\Lambda^{\mathrm T}\Lambda+\mathrm i\gamma'h\) for the matrix of this form on the contour; then \((\epsilon+\mathrm i\gamma')h=M_\gamma-2\epsilon\ell\ell^{\mathrm T}\). Let \(\gamma\le0\) and \(\epsilon\ge0\), not both zero, and write \(\epsilon+\mathrm i\gamma'=|\epsilon+\mathrm i\gamma'|\mathrm e^{\mathrm i\varphi}\) with \(-\pi/2\le\varphi\le0\). Write \(h\in-\mathfrak C_0\) as \(-T^{\mathrm T}\operatorname{diag}(d_\mu)T\) with the arguments \(\beta_\mu\) of \(\mathfrak C_0\). By (ii) of the preceding paragraph at least eight of the \(\beta_\mu\) lie in \((0,\pi/2)\), so \(\operatorname{Re}(\mathrm e^{\mathrm i\varphi}h)\) has at least eight negative eigenvalues, while by the identity above it has at most one when \(\operatorname{Re}M_\gamma\) is positive semidefinite. So \(h\notin-\mathfrak C_0\) persists. For fixed \(\gamma>0\) it fails. The contour \(A_0=\mathrm e^{\mathrm i\pi/8}X_0\), \(A_i=\mathrm e^{-3\pi\mathrm i/8}X_i\) has \(h=-\mathrm e^{\mathrm i\pi/4}\mathbb 1\in-\mathfrak C_0\) and \(\operatorname{Re}\mathrm iS_{\mathrm b}=-S_{\mathrm E}\), and its quadratic exponent has negative real part whenever \(\gamma'>\epsilon\).

(ii) \emph{Necessity at \(N=2\).} At \(N=2\) the condition on \(\operatorname{Re}M\) is necessary for absolute convergence. Suppose \(\operatorname{Re}(a^{\mathrm T}Ma)<0\) for a real \(a\). Take the colour columns \(x^1=Ra'\), \(x^2=z_2/R\) and \(x^3=z_3/R\), with \(a'\) in a small cone around \(a\), \(|z_2|,|z_3|<1\) and \(R\) large. Then \(\det\mathcal G=R^{-2}\det(Z^{\mathrm T}hZ)\) with \(Z=(a',z_2,z_3)\), \(S_{\mathrm b}\) is bounded, and the regulator factor has modulus at least \(\mathrm e^{c\epsilon R^2-C}\). Since \(\det(Z^{\mathrm T}hZ)\) is a nonzero polynomial, the modulus of the integrand of \(I_\epsilon[1]\) integrates to at least a constant times \(\int^\infty R^{-19}\mathrm e^{c\epsilon R^2}\,\mathrm dR=\infty\). So at \(N=2\) every linear contour on which \(I_\epsilon[1]\) converges absolutely has \(\operatorname{Re}M\) positive semidefinite, and the same argument applies to \(\Lambda^{\mathrm T}P\Lambda\) and to \(M_\gamma\). By (i) its \(h\) is not in \(-\mathfrak C_0\), and if the contour is damped, \(h\in\mathfrak C\), then \(h\in\mathfrak C_0\). The bounded no-go of (iii) of the preceding paragraph therefore covers every damped linear contour on which the regulated integrand is absolutely integrable, whatever path leads to it. In particular, no path along which \(\operatorname{Re}\mathrm iS_{\mathrm b}\le0\) while \(h\) leaves the closure of \(\mathfrak C\) ends on a contour with \(h\in-\mathfrak C_0\) on which the regulated integrand is absolutely integrable. For \(N\ge3\) the same holds for the contours on which the regulator factor is bounded by \(1\), the class used in the Stokes arguments above. Non-integrability outside that class is not shown for \(N\ge3\).

(iii) \emph{General positive definite regulators at \(N=2\).} For a general positive definite \(P\) the inertia in (i) is not fixed. Take \(P=\operatorname{diag}(1,1,\tfrac14,1,\dots,1)\) and \(\Lambda=\operatorname{diag}(1,1,2,1,\dots,1)\Lambda'\), where \(\Lambda'\) has rows \((0,0,1)\), \((1,0,0)\) and \((4\mathrm i/5,1,0)\) in the first three indices and is the identity in the others. Then \(\operatorname{Re}\Lambda^{\mathrm T}P\Lambda=\operatorname{diag}(\tfrac9{25},1,1,\dots,1)\), while \(\operatorname{Re}h=\operatorname{diag}(-\tfrac{39}{25},4,-1,1,\dots,1)\) has two negative eigenvalues. Whether such \(P\) admit contours with \(h\in-\mathfrak C_0\) is not decided here. A statement about paths holds instead. At \(N=2\) and for \(h\in\mathfrak C_0\), let \(\langle\operatorname{Tr}X_\alpha X_\beta\rangle_h\) be the ratio of the integrals over real \(X\) of \((\det\mathcal G)^4\mathrm e^{\mathrm iS_{\mathrm b}}\operatorname{Tr}X_\alpha X_\beta\) and of \((\det\mathcal G)^4\mathrm e^{\mathrm iS_{\mathrm b}}\); the denominator does not vanish by (iv) of the preceding paragraph. This ratio equals \(c_{\mathrm E}\mathrm e^{\mathrm i\pi/4}(h^{-1})_{\alpha\beta}\), with \(c_{\mathrm E}=\langle\operatorname{Tr}\tilde A_0^2\rangle_{\mathrm E}>0\). Both sides are holomorphic on \(\mathfrak C_0\), transform alike under real congruences, and agree at \(h=\mathrm e^{\mathrm i\pi/4}\mathbb 1\). So they agree at every \(\mathrm e^{\mathrm i\pi/4}T^{\mathrm T}T\), and by the identity theorem, as in (iv) of the preceding paragraph, everywhere on \(\mathfrak C_0\). On a contour with \(h\in\mathfrak C_0\) this gives \(\langle\operatorname{Tr}A_\mu A_\nu\rangle=c_{\mathrm E}\mathrm e^{\mathrm i\pi/4}\eta_{\mu\nu}\), which is the limit of the paragraph on the regulated Lorentzian integral. The integrand is even in \(h\), so on a contour with \(h\in-\mathfrak C_0\) the same ratio gives \(-c_{\mathrm E}\mathrm e^{\mathrm i\pi/4}\eta_{\mu\nu}\). Now let \(\Lambda(\tau)\), \(0\le\tau\le1\), be piecewise continuously differentiable, with \(\Lambda(0)\) real, \(\operatorname{Re}\mathrm iS_{\mathrm b}\le0\) on every contour of the path, \(\operatorname{Re}\Lambda(\tau)^{\mathrm T}P\Lambda(\tau)\) positive definite and \(h(1)\in\mathfrak C\). Stokes' theorem at fixed \(\epsilon\), as for the \(\mathcal R_u\), and dominated convergence on the last contour give the limit of \(\langle\operatorname{Tr}A_\mu A_\nu\rangle_\epsilon\) as its value on that contour. That paragraph gives this limit as \(c_{\mathrm E}\mathrm e^{\mathrm i\pi/4}\eta_{\mu\nu}\) for every positive definite \(P\). Hence \(h(1)\in\mathfrak C_0\). For every such regulator, no path of linear contours along which \(\operatorname{Re}\mathrm iS_{\mathrm b}\le0\) and the regulator decays joins the Hermitian contour to \(-\mathfrak C_0\).

(iv) \emph{What the contours of \(-\mathfrak C_0\) compute.} For \(h\in-\mathfrak C_0\) write \(\Lambda=\mathrm i\Lambda'\), so that the contour \(\Lambda'\) has \(-h\in\mathfrak C_0\). The functions \(S_{\mathrm b}\) and \(\operatorname{Pf}\mathcal M\) are homogeneous of degrees \(4\) and \(8(N^2-1)\), so the weight at \(\epsilon=0\) is the same on both contours up to a constant factor. The expectation value of \(O\) on the contour \(\Lambda\) is therefore that of \(A\mapsto O(\mathrm iA)\) on \(\Lambda'\), for every \(N\) where these integrals converge absolutely. For \(O=\operatorname{Tr}\mathrm e^{\mathrm ik^\mu A_\mu}\) this is the Fourier mode at the momentum \(\mathrm ik\) on a contour of \(\mathfrak C_0\). At \(N=2\) and real timelike \(k\), the value \(F_{\mathrm E}(\mathrm e^{\mathrm i\pi/4}\kappa^2)\) of the preceding paragraph is the number that (v) there gives as the real spacelike limit at \(k^\mu k_\mu=\kappa^2\). No contour with \(\operatorname{Re}M\) positive definite makes the mode bounded at an imaginary timelike momentum \(k^\mu=\mathrm i\kappa\delta^\mu_0\). For real \(q=\Lambda^{\mathrm T}k\) one has \(\sum_\mu(k^\mu)^2=q^{\mathrm T}M^{-1}q\), whose real part is positive because \(\operatorname{Re}M^{-1}\) is positive definite, while \(\sum_\mu(k^\mu)^2=-\kappa^2\) at that momentum. The value on \(-\mathfrak C_0\) at real timelike \(k\) is therefore not obtained in this way as a limit of the regulated integral, at real or at imaginary timelike momentum.

So, for linear contours, Definition~(A) with \(\gamma\le0\) or a regulator as in (i) with \(\gamma=0\), and \(\epsilon\to0\) taken at fixed \(N\), the damped contours on which the regulated integrand is absolutely integrable (at \(N=2\)) or the regulator factor is bounded by \(1\) (at \(N\ge3\)) all lie in \(\mathfrak C_0\), the class of the Hermitian contour. None of them reaches a real causal momentum. For any positive definite regulator, the same holds at \(N=2\) for the end points of paths as in (iii). The contours of \(-\mathfrak C_0\) that make the mode bounded at real timelike \(k\) compute observables at imaginary momenta. This is no statement about non-linear contours, about linear maps that also mix the colour components, about contours with \(h\in-\mathfrak C_0\) for general positive definite \(P\), about fixed \(\gamma>0\), about whether \(F_{\mathrm E}\) continues beyond \(\arg t=-3\pi/4\), about the pointwise limit at real timelike or null \(k\), about \(N\to\infty\) taken first, or about Definitions~(B) and~(C). The inertia found in (i) is that of the \(\eta\) put into \(S_{\mathrm b}\), which the regulator carries over to every admissible contour but does not select, so Assumption~\ref{ass:lorentzian} is still required.
